\pagebreak

\section{Spartan6, Virtex4, Virtex5, Virtex6 and series 7 DSP48 block Library File - xilinx/dsp48.3pl}

    \index{library!xilinx/dsp48.3pl}
    This library creates a DSP48 block for Spartan6, Virtex4, Virtex5, Virtex6
    and series 7 FPGAs.
    
    This library must be incorporated by using a {\bf module} statement -
    \begin{lstlisting}[gobble=4, language=threepl]
       module "xilinx/dsp48.3pl"
    \end{lstlisting}



    \subsection{library modules}

        \subsubsection{dspaddsub - DSP add/subtract}\label{sec:lmdspaddsub}

            {\bf Library: xilinx/dsp48.3pl}

            {\bf Prototype:} \\
            \vsp
            \begin{lstlisting}[gobble=12, language=threepl]
               global module
               dspaddsub (
                   value(out),
                   value(pat, "log")
                   <-
                   value(a),
                   value(b),
                   value(add, "log", true),
                   value(gate, "log", true),
                   int(areg, breg, preg),
                   int(pattern),
                   int(mask, 0xffffffffffff)
               ) {}
            \end{lstlisting}

            {\bf Output parameters:} \\
            {\bf out} result.\\
            {\bf pat} pattern matched.\\


            {\bf Input parameters:} \\
            {\bf a} first operand.\\
            {\bf b} second operand.\\
            {\bf add} add or subtract.\\
            {\bf gate} gate second operand.\\
            {\bf areg} number of registers in A input.\\
            {\bf breg} number of registers in B input.\\
            {\bf preg} number of registers on DSP48E output.\\
            {\bf pattern} pattern to match against result.\\
            {\bf mask} mask to select pattern bits.\\

            {\bf description:} \\
            This module uses the DSP48 to add or subtract two operands. The
            operands may be immediate, input, value, static or queue mode and may
            be of type "uint:", "int:", "ufixed:", or "fixed:". The types may be
            mixed. The operands may be up to 48 bits signed or 47 bits unsigned
            and need not be the same size. The output result type is given by
            the following table -

            \begin{tabular}{| l || r | r | r | r |}
            \hline
                    & uint:   & int:   & ufixed: & fixed: \\
            \hline
            \hline
            uint:   & uint:   & int:   & ufixed: & fixed: \\
            \hline
            int:    & int:    & int:   & fixed:  & fixed: \\
            \hline
            ufixed: & ufixed: & fixed: & ufixed: & fixed: \\
            \hline
            fixed:  & fixed:  & fixed: & fixed:  & fixed: \\
            \hline
            \end{tabular}

            The output result binary point offset, i.e. number of fraction
            bits, is equal to the largest binary point offset of the two
            operands. The output result integer part width is equal to the
            largest integer part width of the two operands plus one. The
            full width of the result cannot exceed 48 bits. The result will
            be cast to fit the type of the argument passed to the first
            output parameter {\bf out}.

            The argument supplied to the first output parameter, {\bf out}, is a
            value mode variable which should have the type "empty" to allow it
            to take the result type.

            The third input argument {\bf add} is type "log" and is true for
            addition and false for subtraction. If omitted, the operation is
            addition. An immediate mode argument may be supplied for a
            fixed operator or a target value may be used to select the
            operator dynamically.

            The fourth input argument {\bf gate} is type "log" and is true for
            addition or subtraction of the second operand and false to pass the
            first operand, i.e. add or subtract 0. If omitted, the operation is
            addition/subtraction. An immediate mode argument may also be
            supplied.

            The fifth and sixth input parameters {\bf areg} and {\bf breg} 
            are type "int" and specify the number of registers to insert
            in the {\bf a} and {\bf b} inputs respectively. Allowed
            values for {\bf areg} are 0 (default), 1 and 2 and for
            {\bf breg} 0 (default), and 1.

            The seventh input parameter{\bf preg} specifies the number of
            registers on the result and pattern match outputs. Allowed values
            are 0 (default) or 1. If {\bf preg} is 0 and the outputs are
            assigned to static variables those variables may be subsumed into
            the DSP48 output register.\footnote{this will only be done if the
            static variables are not initialised (0 is allowed), are not
            assigned any other value but a DSP48 output (can be assigned 0/false
            or reset) and if both outputs are used are assigned (or reset) in
            the same way (have a common clock-enable and reset). This is because
            the DSP48 output registers for the arithmetic output and
            the pattern match output share a clock enable and reset and they
            cannot be initialised). An assignment of 0/false to these static
            variables becomes a reset since by default it is the assignment of
            the initial values. Any other assignments would require an input
            multiplexer between the DSP48 output and the registers.}

            The eigth parameter {\bf pattern} is an optional bit pattern to
            match against the arithmetic result. If they are identical the
            second output argument {\bf pat} will be true. The last parameter
            {\bf mask} selects which bits from ANDing the result and the pattern
            are to be used. All bits passed by the mask must be true for the
            match result to be true. If {\bf mask} is not passed an argument
            then all bits are matched to the pattern.

            The second output parameter {\bf pat} provides the result of the
            pattern match, if used.

            If no internal registers are used the input arguments may contain
            queue reads, in which case the arguments to output parameters {\bf
            out} and {\bf pat} will be available for reading when all input
            arguments are available. If any inserted registers are specified
            ({\bf areg}, {\bf breg}, {\bf creg} or {\bf preg} are
            non-zero) the inputs cannot contain queue reads and the code will
            exit with a fatal error.




        \subsubsection{dspcounter - DSP counter}\label{sec:lmdspcounter}

            {\bf Library: xilinx/dsp48.3pl}

            {\bf Prototype:} \\
            \vsp
            \begin{lstlisting}[gobble=12, language=threepl]
               global module
               dspcounter (
                   value(out),
                   value(pat, "log")
                   <-
                   value(step, "log", true),
                   value(reset, "log", false),
                   value(pattern, "bits:48", 0),
                   int(mask, 0xffffffffffff),
                   value(up, "log", true)
               ) {}
            \end{lstlisting}

            {\bf Output parameters:} \\
            {\bf out} counter value.\\
            {\bf pat} pattern matched.\\


            {\bf Input parameters:} \\
            {\bf step} advance the count.\\
            {\bf reset} clear the count.\\
            {\bf pattern} pattern to match against count.\\
            {\bf mask} mask to select pattern bits.\\
            {\bf up} up or down.\\

            {\bf description:} \\
            This module uses the DSP48 as a 48 bit counter. The count will
            be cast to fit the type of the argument passed to the first
            output parameter {\bf out}.

            The first input parameter, {\bf step}, is a boolean which will
            advance the count when true. If it is not passed an argument
            then by default the counter will advance on every clock cycle.

            The second input parameter, {\bf reset}, is a boolean which will
            synchronously reset the counter to zero when true.

            The third parameter {\bf pattern} is an optional bit pattern to
            match against the count. If they are identical the second output
            argument {\bf pat} will be true. The pattern may be either immediate
            mode (constant) or target mode to compare against a target variable
            or expression. The fourth parameter {\bf mask} selects which bits
            from ANDing the count and the pattern are to be used. All bits
            passed by the mask must be true for the match result to be true. If
            {\bf mask} is not passed an argument then all bits are matched to
            the pattern.

            The last input parameter determines the count direction, true for
            up and false for down. It may of course have an immediate or
            target argument depending on whether the count direction is to
            be changed dynamically or fixed. If the argument is omitted
            the direction is up.

            The second output parameter {\bf pat} provides the result of the
            pattern match, if used.

            The input arguments must not contain queue reads.




        \subsubsection{dsplog - DSP logical operation}\label{sec:lmdsplog}

            {\bf Library: xilinx/dsp48.3pl}

            {\bf Prototype:} \\
            \vsp
            \begin{lstlisting}[gobble=12, language=threepl]
               global module
               dsplog (
                   value(out),
                   value(pat, "log")
                   <-
                   value(a),
                   value(b),
                   value(op, "uint:3"),
                   int(pattern),
                   int(mask, 0xffffffffffff)
               ) {}
            \end{lstlisting}

            {\bf Output parameters:} \\
            {\bf out} counter value.\\
            {\bf pat} pattern matched.\\


            {\bf Input parameters:} \\
            {\bf a} first input operand.\\
            {\bf b} second input operand.\\
            {\bf op} operator.\\
            {\bf pattern} pattern to match against result.\\
            {\bf mask} mask to select pattern bits.\\

            {\bf description:} \\
            This module uses the DSP48 as a 48 bit counter. The result
            is a logical operation on the two input operands.

            Input parameters {\bf a} and {\bf b} are the operands and
            {\bf op} is an operation code. The operation codes are -

            \begin{tabular}{| r r r |}
            \hline
            4   & 0b000100    & \verb'  a ^ b ' \\
            5   & 0b000101    & \verb'~(a ^ b)' \\
            6   & 0b000110    & \verb'~(a ^ b)' \\
            7   & 0b000111    & \verb'  a ^ b ' \\
            12  & 0b001100    & \verb'  a & b ' \\
            13  & 0b001010    & \verb'  a & ~b' \\
            14  & 0b001110    & \verb'~(a & b)' \\
            15  & 0b001111    & \verb' ~a | b ' \\
            36  & 0b100100    & \verb'~(a ^ b)' \\
            37  & 0b100101    & \verb'  a ^ b ' \\
            38  & 0b100110    & \verb'  a ^ b ' \\
            39  & 0b100111    & \verb'~(a ^ b)' \\
            44  & 0b101100    & \verb'  a | b ' \\
            45  & 0b101101    & \verb'  a | ~b' \\
            46  & 0b101110    & \verb'~(a ^ b)' \\
            47  & 0b101111    & \verb' ~a & b ' \\
            \hline
            \end{tabular}

            The argument to {\bf op} may be immediate mode for a fixed
            operation or may be target mode to dynamically select the
            operation.

            The operands can be any width up to 48 bits and can be any type.

            The fourth parameter {\bf pattern} is an optional bit pattern to
            match against the result. If they are identical the second
            output argument {\bf pat} will be true. The last parameter
            {\bf mask} selects which bits from ANDing the result and the pattern
            are to be used. All bits passed by the mask must be true for the
            match result to be true. If {\bf mask} is not passed an argument
            then all bits are matched to the pattern.

            The second output parameter {\bf pat} provides the result of the
            pattern match, if used.

            The input arguments may contain queue reads in which case the
            arguments to output parameters {\bf out} and {\bf pat} will
            be available for reading when all input arguments are
            available.




        \subsubsection{dspmul - DSP multiply}\label{sec:lmdspmul}

            {\bf Library: xilinx/dsp48.3pl}

            {\bf Prototypes:} \\
            \vsp
            \begin{lstlisting}[gobble=12, language=threepl]
               global module
               dspmul (
                   value(out), value(pat, "log")
                   <-
                   value(a),
                   value(b),
                   int(areg, breg, mreg, preg),
                   int(pattern),
                   int(mask, 0xffffffffffff)
               ) {}
            \end{lstlisting}

            {\bf Output parameters:} \\
            {\bf out} result.\\
            {\bf pat} pattern matched.\\


            {\bf Input parameters:} \\
            {\bf a} first operand.\\
            {\bf b} second operand.\\
            {\bf areg} number of registers in A input.\\
            {\bf breg} number of registers in B input.\\
            {\bf mreg} number of registers on multiplier output.\\
            {\bf preg} number of registers on DSP48 output.\\
            {\bf pattern} pattern to match against result.\\
            {\bf mask} mask to select pattern bits.\\

            {\bf description:} \\
            These modules use the DSP48 to multiply two signed operands. The
            operands may be immediate, input, value, static or queue mode and may
            be of type "uint:", "int:", "ufixed:", or "fixed:". The types may be
            mixed. The result type is determined by the input operand type
            nearest the top of the following table -

            \begin{tabular}{| l |}
            \hline
            fixed   \\ \hline
            ufixed  \\ \hline
            int     \\ \hline
            uint    \\ \hline
            \end{tabular}

            The first operand, {\bf a}, is 25 bits maximum (24 bits if unsigned)
            and the second operand, {\bf b}, is 18 bits maximum (17 bits if
            unsigned).

            The output result binary point offset, i.e. number of fraction bits,
            is equal to the sum of the offsets of the operands. The product
            integer part width is equal to the sum of the integer widths of its
            two operands. The result will be cast to fit the type of the
            argument passed to the first output parameter {\bf out}.

            The argument supplied to the first output parameter, {\bf out}, is a
            value mode variable which may have the type "empty" to allow it
            to take the result type without casting.

            The third and fourth input parameters {\bf areg} and {\bf breg}
            are type "int" and specify the number of registers to insert
            in the {\bf a} and {\bf b} inputs respectively. Allowed
            values for {\bf areg} and {\bf breg} are 0 (default), 1 and 2.

            The fifth input parameter {\bf mreg} is type "int" and specifies the
            number of registers to insert between the multiplier and the
            adder/subtracter which in this case simply passes the product.
            Allowed values are 0 (default) or 1.

            The sixth input parameter{\bf preg} specifies the number of
            registers on the result and pattern match outputs. Allowed values
            are 0 (default) or 1. If {\bf preg} is 0 and the outputs are
            assigned to static variables those variables may be subsumed into
            the DSP48 output register.\footnote{this will only be done if the
            static variables are not initialised (0 is allowed), are not
            assigned any other value but a DSP48 output (can be assigned 0/false
            or reset) and if both outputs are used are assigned (or reset) in
            the same way (have a common clock-enable and reset). This is because
            the DSP48 output registers for the arithmetic output and
            the pattern match output share a clock enable and reset and they
            cannot be initialised). An assignment of 0/false to these static
            variables becomes a reset since by default it is the assignment of
            the initial values. Any other assignments would require an input
            multiplexer between the DSP48 output and the registers.}

            The seventh parameter {\bf pattern} is an optional bit pattern to
            match against the arithmetic result. If they are identical the
            second output argument {\bf pat} will be true. The last input
            parameter {\bf mask} selects which bits from ANDing the result and
            the pattern are to be used. All bits passed by the mask must be true
            for the match result to be true. If {\bf mask} is not passed an
            argument then all bits are matched to the pattern.

            The second output parameter {\bf pat} provides the result of the
            pattern match, if used.

            If no internal registers are used the input arguments may contain
            queue reads, in which case the arguments to output parameters {\bf
            out} and {\bf pat} will be available for reading when all input
            arguments are available. If any inserted registers are specified
            ({\bf areg}, {\bf breg}, {\bf mreg} or {\bf preg} are non-zero) the
            inputs cannot contain queue reads and the code will exit with a fatal
            error.




        \subsubsection{dspmuladd - DSP multiply and add}\label{sec:lmdspmuladd}
        \subsubsection{dspmulsub - DSP multiply and subtract}\label{sec:lmdspmulsub}
        \subsubsection{dspsubmul - DSP subtract and multiply}\label{sec:lmdspsubmul}

            {\bf Library: xilinx/dsp48.3pl}

            {\bf Prototypes:} \\
            \vsp
            \begin{lstlisting}[gobble=12, language=threepl]
               global module
               dspmuladd (
                   value(out), value(pat, "log")
                   <-
                   value(a),
                   value(b),
                   value(c),
                   int(areg, breg, creg, mreg, preg),
                   int(pattern),
                   int(mask, 0xffffffffffff)
               ) {}
            \end{lstlisting}
            \vsp
            \begin{lstlisting}[gobble=12, language=threepl]
               global module
               dspmulsub (
                   value(out), value(pat, "log")
                   <-
                   value(a),
                   value(b),
                   value(c),
                   int(areg, breg, creg, mreg, preg),
                   int(pattern),
                   int(mask, 0xffffffffffff)
               ) {}
            \end{lstlisting}
            \vsp
            \begin{lstlisting}[gobble=12, language=threepl]
               global module
               dspsubmul (
                   value(out), value(pat, "log")
                   <-
                   value(a),
                   value(b),
                   value(c),
                   int(areg, breg, creg, mreg, preg),
                   int(pattern),
                   int(mask, 0xffffffffffff)
               ) {}
            \end{lstlisting}

            {\bf Output parameters:} \\
            {\bf out} result.\\
            {\bf pat} pattern matched.\\


            {\bf Input parameters:} \\
            {\bf a} first operand.\\
            {\bf b} second operand.\\
            {\bf c} third operand.\\
            {\bf areg} number of registers in A input.\\
            {\bf breg} number of registers in B input.\\
            {\bf creg} number of registers in A input.\\
            {\bf mreg} number of registers on multiplier output.\\
            {\bf preg} number of registers on DSP48 output.\\
            {\bf pattern} pattern to match against result.\\
            {\bf mask} mask to select pattern bits.\\

            {\bf description:} \\
            These modules use the DSP48 to combine three operands. The
            operands may be immediate, input, value, static or queue mode and may
            be of type "uint:", "int:", "ufixed:", or "fixed:". The types may be
            mixed. The result type is determined by the input operand type
            nearest the top of the following table -

            \begin{tabular}{| l |}
            \hline
            fixed   \\ \hline
            ufixed  \\ \hline
            int     \\ \hline
            uint    \\ \hline
            \end{tabular}

            The operations performed are -

            \begin{tabular}{| l l |}
            \hline
            dspmuladd & $ a * b + c$ \\
            \hline
            a & 25 bits signed \\
            b & 18 bits signed \\
            c & 48 bits signed \\
            \hline
            \end{tabular}

            \begin{tabular}{| l l |}
            \hline
            dspmulsub & $ a * b - c $ \\
            \hline
            a & 25 bits signed \\
            b & 18 bits signed \\
            c & 48 bits signed \\
            \hline
            \end{tabular}

            \begin{tabular}{| l l |}
            \hline
            dspsubmul & $ a - b * c $ \\
            \hline
            a & 48 bits signed \\
            b & 25 bits signed \\
            c & 18 bits signed \\
            \hline
            \end{tabular}

            Unsigned arguments must be 1 bit smaller.

            The output result binary point offset, i.e. number of fraction bits,
            is equal to the sum of the offsets of the operands to the
            multiplier; if the remaining operand has a larger offset it is
            trimmed at the least significant end. This differs from the way 3PL
            normally handles expressions where significant bits are retained.
            The product integer part width is equal to the sum of the integer
            widths of its two operands. The output result integer part width is
            equal to the largest integer part width of the product and the
            remaining operand plus one. The full width of the result cannot
            exceed 48 bits. The result will be cast to fit the type of the
            argument passed to the first output parameter {\bf out}.

            The argument supplied to the first output parameter, {\bf out}, is a
            value mode variable which may have the type "empty" to allow it
            to take the result type without casting.

            The fourth to sixth input parameters {\bf areg}, {\bf breg} and {\bf
            creg} are type "int" and specify the number of registers to insert
            in the {\bf a}, {\bf b} and {\bf c} inputs respectively. Allowed
            values for {\bf areg} and {\bf breg} are 0 (default), 1 and 2.
            {\bf creg} may be 0 (default) or 1.

            The seventh input parameter {\bf mreg} is type "int" and specifies
            the number of registers to insert between the multiplier and the
            adder/subtracter. Allowed values are 0 (default) or 1.

            The eighth input parameter{\bf preg} specifies the number of
            registers on the result and pattern match outputs. Allowed values
            are 0 (default) or 1. If {\bf preg} is 0 and the outputs are
            assigned to static variables those variables may be subsumed into
            the DSP48 output register.\footnote{this will only be done if the
            static variables are not initialised (0 is allowed), are not
            assigned any other value but a DSP48 output (can be assigned 0/false
            or reset) and if both outputs are used are assigned (or reset) in
            the same way (have a common clock-enable and reset). This is because
            the DSP48 output registers for the arithmetic output and
            the pattern match output share a clock enable and reset and they
            cannot be initialised). An assignment of 0/false to these static
            variables becomes a reset since by default it is the assignment of
            the initial values. Any other assignments would require an input
            multiplexer between the DSP48 output and the registers.}

            The ninth parameter {\bf pattern} is an optional bit pattern to
            match against the arithmetic result. If they are identical the
            second output argument {\bf pat} will be true. The last input
            parameter {\bf mask} selects which bits from ANDing the result and
            the pattern are to be used. All bits passed by the mask must be true
            for the match result to be true. If {\bf mask} is not passed an
            argument then all bits are matched to the pattern.

            The second output parameter {\bf pat} provides the result of the
            pattern match, if used.

            If no internal registers are used the input arguments may contain
            queue reads, in which case the arguments to output parameters {\bf
            out} and {\bf pat} will be available for reading when all input
            arguments are available. If any inserted registers are specified
            ({\bf areg}, {\bf breg}, {\bf creg}, {\bf mreg} or {\bf preg} are
            non-zero) the inputs cannot contain queue reads and the code will
            exit with a fatal error.


    \subsection{library procedures}

        \subsubsection{dsp48 - create a DSP48 block for Virtex4 FPGAs}\label{sec:lpdsp48}

            {\bf Library: xilinx/dsp48.3pl}
            
            {\bf Prototype:} \\
            \vsp
            \begin{lstlisting}[gobble=12, language=threepl]
               global procedure
               dsp48 (
                   value(bcout, "bits:18"),
                   value(pcout, "bits:48"),
                   value(p, "bits:48")
                   <-   
                   str({loc, b_input}),
                   int({areg, breg, creg, mreg, preg, opmodereg,
                        carryinreg, subtractreg, carryinselreg}),

                   clock(clk),
                   value(a, "bits:18"),
                   value(b, "bits:18"),
                   value(c, "bits:48"),
                   value(opmode, "uint:7"),
                   value(subtract, "log"),
                   value(carryinsel, "uint:2"),
                   value(bcin, "bits:18"),
                   value(pcin, "bits:48"),
                   value({carryin, cea, ceb, cec, cem, cep,
                          cectrl, ceinsub, cecarryin, rsta, rstb, rstc, rstm,
                          rstp, rstctrl, rstcarryin}, "log")
               ) {}
            \end{lstlisting}

            {\bf Output parameters:} \\
            {\bf bcout} - 18 bit B carry out   \\
            {\bf pcout} - 48 bit P carry out  \\
            {\bf p} - 48 bit P out    \\

            {\bf Input parameters:} \\
            {\bf AttributesS -} \\
            {\bf loc} - optional location. \\
            {\bf b\_input} - select B input cascade. \\
            {\bf areg} - 0, 1 or 2 for number of pipeline registers. \\
            {\bf breg} - 0, 1 or 2 for number of pipeline registers. \\
            {\bf creg} - 0 or 1 for number of pipeline registers. \\
            {\bf mreg} - 0 or 1 for number of pipeline registers. \\
            {\bf preg} - 0 or 1 for number of pipeline registers. \\
            {\bf opmodereg} - 0 or 1 for number of pipeline registers. \\
            {\bf carryinreg} - 0 or 1 for number of pipeline registers. \\
            {\bf subtractreg} - 0 or 1 for number of pipeline registers. \\
            {\bf carryinselreg} - 0 or 1 for number of pipeline registers. \\
            {\bf Input signals -} \\
            {\bf clk} - clock input. \\
            {\bf a} - 18 bit a input.\\
            {\bf b} - 18 bit b input.\\
            {\bf c} - 48 bit c input.\\
            {\bf opmode} - 7 bit operation mode. \\
            {\bf subtract} - select add or subtract. \\
            {\bf carryinsel} - select carry source. \\
            {\bf bcin} - 18 bit B carry in. \\
            {\bf pcin} - 48 bit P carry in. \\
            {\bf carryin} - carry input. \\
            {\bf cea} - A register clock enable. \\
            {\bf ceb} - B register clock enable. \\
            {\bf cec} - C register clock enable. \\
            {\bf cem} - M register clock enable. \\
            {\bf cep} - P register clock enable. \\
            {\bf cectrl} - opmodereg and carryinreg clock enable. \\
            {\bf ceinsub} - subtract input register clock enable. \\
            {\bf cecarryin} - carry input register clock enable. \\
            {\bf rsta} - reset A register. \\
            {\bf rstb} - reset B register. \\
            {\bf rstc} - reset C register. \\
            {\bf rstm} - reset M register. \\
            {\bf rstp} - reset P register. \\
            {\bf rstctrl} - reset control register. \\
            {\bf rstcarryin} - reset carryin register. \\

            {\bf In-line target execution:} \\
            no

            {\bf description:} \\
            This procedure creates a Virtex4 DSP48 block. The first 11 input
            parameters are attributes, the remaining inputs being input signals.



        \subsubsection{dsp48a - create a DSP48A block for Spartan6 FPGAs}\label{sec:lpdsp48a}

            {\bf Library: xilinx/dsp48.3pl}

            {\bf Prototype:} \\
            \vsp
            \begin{lstlisting}[gobble=12, language=threepl]
               global procedure
               dsp48a (
                   value(bcout, "bits:18"),
                   value(pcout, "bits:48"),
                   value(m, "bits:36"),
                   value(p, "bits:48"),
                   value(carryout, "bits:1"),
                   value(carryoutf, "bits:1")
                   <-   
                   str(loc),
                   str(carryinsel, "CARRYIN"),
                   str(rsttype, "SYNC"),
                   int({a0reg, b0reg, }, 0),
                   int({a1reg, b1reg, creg, dreg, mreg, preg, opmodereg,
                        carryinreg, carryoutreg}, 1)

                   clock(clk),
                   value(a, "bits:18"),
                   value(b, "bits:18"),
                   value(c, "bits:48"),
                   value(d, "bits:18"),
                   value(opmode, "uint:7"),
                   value(pcin, "bits:48"),
                   value(carryin, "bits:1"),
                   value({cea, ceb, cec, ced, cem, cep, ceopmode,
                          cecarryin, rsta, rstb, rstc, rstd, rstm,
                          rstp, rstopmode, rstcarryin}, "log")
               ) {}
            \end{lstlisting}

            {\bf Output parameters:} \\

            {\bf Input parameters:} \\

            {\bf In-line target execution:} \\
            no

            {\bf description:} \\
            This procedure creates a Spartan6 DSP48A block. The first 14 input
            parameters are attributes, the remaining inputs being input signals.


        \subsubsection{dsp48e - create a DSP48E block for Virtex5 FPGAs}\label{sec:lpdsp48e}

            {\bf Library: xilinx/dsp48.3pl}

            {\bf Prototype:} \\
            \vsp
            \begin{lstlisting}[gobble=12, language=threepl]
               global procedure
               dsp48e (
                   value(acout, "bits:30"),
                   value(bcout, "bits:18"),
                   value(pcout, "bits:48"),
                   value(p, "bits:48"),
                   value(carryout, "[4]log"),
                   value({carrycascout, multsignout, patterndetect, patternbdetect,
                           overflow, underflow}, "log")
                   <-   
                   str(loc, a_input, b_input, use_mult, use_simd, sel_pattern, sel_mask,
                       sel_rounding_mask, autoreset_pattern_detect_optinv, use_pattern_detect),
                   int(areg, acascreg, breg, bcascreg, creg, mreg, preg, opmodereg,
                       alumodereg, carryinreg, multcarryinreg, carryinselreg,
                       pattern, mask),
                   log(autoreset_pattern_detect),
                   clock(clk),
                   value(a, "bits:30"),
                   value(b, "bits:18"),
                   value(c, "bits:48"),
                   value(opmode, "uint:7"),
                   value(alumode, "uint:4"),
                   value(carryinsel, "uint:3"),
                   value(acin, "bits:30"),
                   value(bcin, "bits:18"),
                   value(pcin, "bits:48"),
                   value({carryin, cea1, cea2, ceb1, ceb2, cec, cem, cep, cealumode,
                          cectrl, cemultcarryin, cecarryin, rsta, rstb, rstc, rstm,
                          rstp, rstctrl, rstallcarryin, rstalumode, carrycascin,
                          multsignin}, "log")
               ) {}
            \end{lstlisting}

            {\bf Output parameters:} \\
            {\bf acout} - 30 bit A carry out   \\
            {\bf bcout} - 18 bit B carry out   \\
            {\bf pcout} - 48 bit P carry out  \\
            {\bf p} - 48 bit P out    \\
            {\bf carryout} -   \\
            {\bf carrycascout} -   \\
            {\bf multsignout} -   \\
            {\bf patterndetect} -   \\
            {\bf patternbdetect} -   \\
            {\bf overflow} -   \\
            {\bf underflow} -   \\

            {\bf Input parameters:} \\
            {\bf AttributesS -} \\
            {\bf loc} - optional location. \\
            {\bf a\_input} - select A input cascade. \\
            {\bf b\_input} - select B input cascade. \\
            {\bf use\_mult} -  \\
            {\bf use\_simd} -  \\
            {\bf sel\_pattern} -  \\
            {\bf sel\_mask} -  \\
            {\bf sel\_rounding\_mask} -  \\
            {\bf autoreset\_pattern\_detect\_optinv} -  \\
            {\bf use\_pattern\_detect     } -  \\       
            {\bf areg} - 0, 1 or 2 for number of pipeline registers. \\
            {\bf acascreg} -  \\
            {\bf breg} - 0, 1 or 2 for number of pipeline registers. \\
            {\bf bcascreg} -  \\
            {\bf creg} - 0 or 1 for number of pipeline registers. \\
            {\bf mreg} - 0 or 1 for number of pipeline registers. \\
            {\bf preg} - 0 or 1 for number of pipeline registers. \\
            {\bf opmodereg} - 0 or 1 for number of pipeline registers. \\
            {\bf alumodereg} -  \\
            {\bf carryinreg} - 0 or 1 for number of pipeline registers. \\
            {\bf multcarryinreg} - 0 or 1 for number of pipeline registers. \\
            {\bf carryinselreg} - 0 or 1 for number of pipeline registers. \\
            {\bf pattern} -  \\
            {\bf mask} -  \\
            {\bf Input signals -} \\
            {\bf clk} - clock input. \\
            {\bf a} - 18 bit a input.\\
            {\bf b} - 18 bit b input.\\
            {\bf c} - 48 bit c input.\\
            {\bf opmode} - 7 bit operation mode. \\
            {\bf alumode} - 7 bit operation mode. \\
            {\bf carryinsel} - select carry source. \\
            {\bf acin} - 30 bit B carry in. \\
            {\bf bcin} - 18 bit B carry in. \\
            {\bf pcin} - 48 bit P carry in. \\
            {\bf carryin} - carry input. \\            
            {\bf cea1} - A register clock enable. \\
            {\bf cea2} - A register clock enable. \\
            {\bf ceb1} - B register clock enable. \\
            {\bf ceb2} - B register clock enable. \\
            {\bf cec} - C register clock enable. \\
            {\bf cem} - M register clock enable. \\
            {\bf cep} - P register clock enable. \\
            {\bf cealumode} -  \\
            {\bf cectrl} -  \\
            {\bf cemultcarryin} -  \\
            {\bf cecarryin} - carry input register clock enable. \\
            {\bf rsta} - reset A register. \\
            {\bf rstb} - reset B register. \\
            {\bf rstc} - reset C register. \\
            {\bf rstm} - reset M register. \\
            {\bf rstp} - reset P register. \\
            {\bf rstctrl} -  \\
            {\bf rstallcarryin} -  \\
            {\bf rstalumode} -  \\
            {\bf carrycascin} -  \\
            {\bf multsignin} -    \\       

            {\bf In-line target execution:} \\
            no

            {\bf description:} \\
            This procedure creates a Virtex5 DSP48E block. The first 24 input
            parameters are attributes, the remaining inputs being input signals.


        \subsubsection{dsp48e1 - create a DSP48E1 block for Virtex6 and series 7 FPGAs}\label{sec:lpdsp48e1}

            {\bf Library: xilinx/dsp48.3pl}

            {\bf Prototype:} \\
            \vsp
            \begin{lstlisting}[gobble=12, language=threepl]
               global procedure
               dsp48e1 (
                   value(acout, "bits:30"),
                   value(bcout, "bits:18"),
                   value(pcout, "bits:48"),
                   value(p, "bits:48"),
                   value(carryout, "[4]log"),
                   value({carrycascout, multsignout, patterndetect, patternbdetect,
                           overflow, underflow}, "log")
                   <-   
                   str(loc, a_input, b_input, use_mult, use_simd, sel_pattern, sel_mask,
                       sel_rounding_mask, autoreset_pattern_detect_optinv, use_pattern_detect),
                   int(areg, acascreg, breg, bcascreg, creg, mreg, preg, opmodereg,
                       alumodereg, carryinreg, multcarryinreg, carryinselreg,
                       pattern, mask),
                   log(autoreset_pattern_detect),
                   clock(clk),
                   value(a, "bits:30"),
                   value(b, "bits:18"),
                   value(c, "bits:48"),
                   value(d, "bits:25"),
                   value(opmode, "uint:7"),
                   value(alumode, "uint:4"),
                   value(carryinsel, "uint:3"),
                   value(acin, "bits:30"),
                   value(bcin, "bits:18"),
                   value(pcin, "bits:48"),
                   value({carryin, cea1, cea2, ceb1, ceb2, cec, cem, cep, cealumode,
                          cectrl, cemultcarryin, cecarryin, rsta, rstb, rstc, rstm,
                          rstp, rstctrl, rstallcarryin, rstalumode, carrycascin,
                          multsignin}, "log")
               ) {}
            \end{lstlisting}

            {\bf Output parameters:} \\
            {\bf acout} - 30 bit A carry out   \\
            {\bf bcout} - 18 bit B carry out   \\
            {\bf pcout} - 48 bit P carry out  \\
            {\bf p} - 48 bit P out    \\
            {\bf carryout} -   \\
            {\bf carrycascout} -   \\
            {\bf multsignout} -   \\
            {\bf patterndetect} -   \\
            {\bf patternbdetect} -   \\
            {\bf overflow} -   \\
            {\bf underflow} -   \\

            {\bf Input parameters:} \\
            {\bf AttributesS -} \\
            {\bf loc} - optional location. \\
            {\bf a\_input} - select A input cascade. \\
            {\bf b\_input} - select B input cascade. \\
            {\bf use\_mult} -  \\
            {\bf use\_simd} -  \\
            {\bf sel\_pattern} -  \\
            {\bf sel\_mask} -  \\
            {\bf sel\_rounding\_mask} -  \\
            {\bf autoreset\_pattern\_detect\_optinv} -  \\
            {\bf use\_pattern\_detect     } -  \\       
            {\bf areg} - 0, 1 or 2 for number of pipeline registers. \\
            {\bf acascreg} -  \\
            {\bf breg} - 0, 1 or 2 for number of pipeline registers. \\
            {\bf bcascreg} -  \\
            {\bf creg} - 0 or 1 for number of pipeline registers. \\
            {\bf mreg} - 0 or 1 for number of pipeline registers. \\
            {\bf preg} - 0 or 1 for number of pipeline registers. \\
            {\bf opmodereg} - 0 or 1 for number of pipeline registers. \\
            {\bf alumodereg} -  \\
            {\bf carryinreg} - 0 or 1 for number of pipeline registers. \\
            {\bf multcarryinreg} - 0 or 1 for number of pipeline registers. \\
            {\bf carryinselreg} - 0 or 1 for number of pipeline registers. \\
            {\bf pattern} -  \\
            {\bf mask} -  \\
            {\bf Input signals -} \\
            {\bf clk} - clock input. \\
            {\bf a} - 18 bit a input.\\
            {\bf b} - 18 bit b input.\\
            {\bf c} - 48 bit c input.\\
            {\bf d} - 25 bit c input.\\
            {\bf opmode} - 7 bit operation mode. \\
            {\bf alumode} - 7 bit operation mode. \\
            {\bf carryinsel} - select carry source. \\
            {\bf acin} - 30 bit B carry in. \\
            {\bf bcin} - 18 bit B carry in. \\
            {\bf pcin} - 48 bit P carry in. \\
            {\bf carryin} - carry input. \\            
            {\bf cea1} - A register clock enable. \\
            {\bf cea2} - A register clock enable. \\
            {\bf ceb1} - B register clock enable. \\
            {\bf ceb2} - B register clock enable. \\
            {\bf cec} - C register clock enable. \\
            {\bf cem} - M register clock enable. \\
            {\bf cep} - P register clock enable. \\
            {\bf cealumode} -  \\
            {\bf cectrl} -  \\
            {\bf cemultcarryin} -  \\
            {\bf cecarryin} - carry input register clock enable. \\
            {\bf rsta} - reset A register. \\
            {\bf rstb} - reset B register. \\
            {\bf rstc} - reset C register. \\
            {\bf rstm} - reset M register. \\
            {\bf rstp} - reset P register. \\
            {\bf rstctrl} -  \\
            {\bf rstallcarryin} -  \\
            {\bf rstalumode} -  \\
            {\bf carrycascin} -  \\
            {\bf multsignin} -    \\       

            {\bf In-line target execution:} \\
            no

            {\bf description:} \\
            This procedure creates a Virtex6 DSP48E1 block. The first 24 input
            parameters are attributes, the remaining inputs being input signals.




    \subsection{library functions}
    
        None.






